% Model Evaluation slides
\begin{frame}{Kết quả đánh giá mô hình}
  \centering
  \begin{tabular}{lcc}
    \toprule
    \textbf{Metric} & \textbf{Giá trị} & \textbf{Ý nghĩa} \\
    \midrule
    \textcolor{accentblue}{\textbf{ROC AUC}} & \Large \textbf{0.8031} & Phân biệt tốt, vượt mức ngẫu nhiên xa \\
    \textcolor{accentblue}{\textbf{Recall}} & \Large \textbf{0.6383} & Tóm được khoản 64\% sự cố \\
    Accuracy & 0.7517 & Độ chính xác chung \\
    Precision & 0.3488 & Hiệu suất trong cảnh báo \\
    F1 & 0.4511 & Cân bằng Precision/Recall \\
    \bottomrule
  \end{tabular}
\end{frame}

\begin{frame}{Đánh giá mô hình - Diễn giải}
  \begin{alertblock}{}
    Mô hình ưu tiên \textbf{Recall} (hạn chế bỏ sót nhân tài) - phù hợp với mục tiêu nhân sự
  \end{alertblock}
  
  \vspace{0.4cm}
  \textbf{Ý nghĩa đánh đổi (Trade-off):}
  \begin{itemize}
    \item Nếu mất nhân sự tốn $\gg$ chi phí một buổi trò chuyện, ta cần giữ threshold thấp (như 0.5) để tối đa Recall.
    \item Mô hình chấp nhận một tỷ lệ nhất định báo động ảo (False Positive) để đánh đổi sự an toàn trong quản trị nhân lực.
  \end{itemize}
\end{frame}

\begin{frame}{ROC Curve}
  \begin{figure}
    \centering
    \includegraphics[width=0.65\textwidth]{roc_curve.png}
  \end{figure}
\end{frame}

\begin{frame}{Confusion Matrix}
  \begin{figure}
    \centering
    \includegraphics[width=0.55\textwidth]{confusion_matrix.png}
  \end{figure}
\end{frame}

\begin{frame}{Confusion Matrix - Diễn giải}
  \textbf{Phát hiện được khoảng 64\% nguy cơ (Recall = 0.638)}
  
  \vspace{0.4cm}
  \textbf{Thành phần:}
  \begin{itemize}
    \item \textbf{True Positive (TP)}: Nhân viên nghỉ việc được phát hiện và HR có thể cứu vãn
    \item \textbf{False Negative (FN)}: Chảy máu chất xám do hệ thống bỏ sót
    \item \textbf{False Positive (FP)}: Tốn kém dư thừa cho nhân viên không có ý định nghỉ
    \item \textbf{True Negative (TN)}: Quản trị phần đông người tại vị an toàn
  \end{itemize}
\end{frame}
